% ===================================================================
%  TAFEL Nr. 14 — Die Straße. --- Die Gewerbetreibenden.
%  Traduction 2026 d'apres texto/io, controlee sur texto/fr et
%  texto/en. Consigne : voir texto/de/10-tafel-01.tex.
%
%  CETTE PAGE POSE A LA CONSIGNE DE 2026 SA QUESTION LA PLUS DURE, ET
%  IL FAUT Y REPONDRE PAR ECRIT.
%
%  L'allemand marque le sexe de celui qui exerce un metier. Cette
%  planche nomme une dizaine de femmes par leur travail, et chacune
%  prend le suffixe : Näherin, Wäscherin,
%  Stenotypistin, Verkäuferin,
%  Blumenhändlerin, Kassiererin,
%  Putzmacherin, Erzieherin, Gesellschafterin,
%  Kinderfrau. Le francais du livret les nomme de meme —
%  couturieres, blanchisseuse, marchande, gouvernante — et l'ido aussi,
%  avec son « -ino ».
%
%  L'allemand de 2026 est en train de defaire ce systeme, et l'on
%  pourrait croire que la consigne oblige a suivre. Elle n'y oblige
%  pas, pour la meme raison qu'au tableau 7 : la regle dit qu'on
%  modernise la LANGUE et non les CHOSES. Or ici les deux formes
%  existent toujours l'une et l'autre en 2026, et le suffixe n'est pas
%  un archaisme — il est la forme normale quand on designe UNE FEMME
%  PARTICULIERE, ce qui est exactement le cas de chacune des dix.
%  LA PLANCHE MONTRE DES PERSONNES, NON DES FONCTIONS. On garde donc
%  le suffixe, et l'on ecrit ici que c'est une decision et non une
%  inadvertance.
%
%  Le tableau 7 avait garde « Knecht » et « Magd » parce que le mot
%  etait mort avec la chose. Ici c'est l'inverse : on garde parce que
%  le mot est vivant. Les deux fois, la question posee est la meme —
%  qu'est-ce qui a change depuis 1926, le mot ou le monde ? — et les
%  deux fois la reponse se lit dans l'usage d'aujourd'hui, non dans
%  celui d'hier.
%
%  ET L'ALINEA 17 PORTE LE PLUS VIEIL EXPORT ALLEMAND DE TOUTE LA
%  SERIE. L'imprimerie du journal emploie un Setzer, un
%  Drucker, un Buchbinder : trois mots allemands
%  partis d'un atelier de Mayence vers 1450 et arrives tels quels ou
%  presque en polonais (drukarz), en tcheque (tiskař, calque), en
%  hongrois (nyomdász, calque), en russe et en yiddish. C'est le
%  meme mouvement que le « Gletscher » du tableau 9 — le mot part avec
%  la chose —, mais quatre siecles plus tot, et c'est le plus grand
%  depart de mots que l'allemand ait jamais fait.
%
%  Le reste de la page se nomme tout seul, par composition : Uhrmacher,
%  Hutmacher, Handschuhmacher, Buchhändler, Möbelhändler,
%  Eisenwarenhändler, Weinhändler, Kohlenhändler, Schirmmacher,
%  Büchsenmacher, Wagenbauer, Zahnarzt, Geldwechsler. Il n'y a que
%  cinq emprunts sur quarante metiers — Juwelier, Friseur, Konditor,
%  Optiker, Fotograf — et ce sont les cinq metiers de luxe ou
%  d'invention recente. C'est la regle du tableau 13 et celle du
%  tableau 5, verifiees ensemble sur une seule rue.
% ===================================================================

%%K t14-apar-1 apar
\VUsaut{4.0mm}
\VUtitre{15.0pt}{60}{TAFEL Nr. 14}
\VUsaut{3.0mm}

%%K t14-tit-1 sub
\VUpk{11.4pt}{40}{Die Straße. --- Die Gewerbetreibenden.}
\VUsaut{3.0mm}

%%K t14-01-1 p
1. --- Mein Freund Johann, der auf dem Land wohnt, ist gekommen, um drei Tage bei
meiner Familie zu verbringen. \VUgras{Bis} jetzt hatte er nie Gelegenheit gehabt,
eine große Stadt zu sehen, und so verbringen wir unsere ganzen Tage mit Spaziergängen,
und ich erkläre ihm alles, was er nicht kennt.

%%K t14-02-1 p
2. --- Zuerst sind wir zur \VUgras{Sternwarte} \textsuperscript{(1)} gegangen, die ihn
sehr interessiert hat. Auf dem Rückweg durch die \VUgras{Straße}
\textsuperscript{(3)}, in der das \VUgras{Dampfbad} \textsuperscript{(2)} liegt, kam
uns ein \VUgras{Leichenzug} \textsuperscript{(4)} entgegen. An der Spitze des
\VUgras{Trauerzuges} ging die \VUgras{Geistlichkeit}, vor dem \VUgras{Leichenwagen},
in den der \VUgras{Sarg} gestellt worden war. Viele Freunde folgten der Familie in
\VUgras{Trauer}.

%%K t14-02-2 p
Als der Zug vorbei war, gingen wir vor der großen \VUgras{Bierhalle}
\textsuperscript{(5)} über die Straße, wo viele \VUgras{Gäste} auf der Terrasse
\VUgras{Bier} tranken, geschützt von der gläsernen \VUgras{Markise}
\textsuperscript{(6)}.

%%K t14-03-1 p
3. --- Johann war das \VUgras{Wappen} zwischen dem Laden des
\VUgras{Weinhändlers} \textsuperscript{(7)} und dem des
\VUgras{Musikalienhändlers} \textsuperscript{(10)} ein Rätsel. Er fragte mich, was
es bedeute; ich sagte ihm, es bezeichne das englische \VUgras{Konsulat}
\textsuperscript{(8)}, und zeigte ihm den \VUgras{Konsul} \textsuperscript{(9)}, der
auf dem Balkon stand.

%%K t14-tit-2 sub
\VUpk{10.2pt}{40}{Schaufenster ansehen.}
\VUsaut{3.0mm}

%%K t14-04-1 p
4. --- Natürlich blieben wir vor der Auslage mit den
\VUgras{Musikinstrumenten} stehen, den \VUgras{Harmonien}, \VUgras{Orgeln} und
\VUgras{Klavieren}; wir sahen uns die bebilderten Umschläge der \VUgras{Noten} und
der \VUgras{Klavierstücke} an und bewunderten die vergoldete hölzerne \VUgras{Leier},
die als Ladenschild dient.

%%K t14-05-1 p
5. --- Die Staubwolken, die die \VUgras{Straßenkehrer} \textsuperscript{(11)} und die
\VUgras{Kehrmaschine} \textsuperscript{(12)} aufwirbelten, zwangen uns, an den schönen
\VUgras{Zimmereinrichtungen} beim \VUgras{Möbelhändler} \textsuperscript{(13)} und an
der Auslage mit \VUgras{Öfen} und \VUgras{Lampen} beim
\VUgras{Eisenwarenhändler} \textsuperscript{(14)} rasch vorbeizugehen.

%%K t14-06-1 p
6. --- An dieser Stelle mussten wir tatsächlich den Bürgersteig verlassen, weil er
von Paketen versperrt war, die man aus einem \VUgras{Möbelwagen}
\textsuperscript{(16)} in einen eben vermieteten \VUgras{Laden}
\textsuperscript{(15)} trug.

%%K t14-06-2 p
Dadurch entgingen uns die schönen Wagen beim \VUgras{Wagenbauer}
\textsuperscript{(17)}, nicht aber der \VUgras{Packer} \textsuperscript{(19)}, dem wir
zusahen, wie er für seinen Nachbarn, den \VUgras{Fahrrad- und Autohändler}
\textsuperscript{(18)}, eine Kiste zimmerte. Dann gingen wir nach Hause, weil es schon
ziemlich spät war.

%%K t14-tit-3 sub
\VUpk{10.2pt}{40}{Ein Gang durch die Stadt.}
\VUsaut{3.0mm}

%%K t14-07-1 p
7. --- Am nächsten Tag schlug meine Mutter vor, Johann fotografieren zu lassen, und
bat mich, mit ihm zum \VUgras{Fotografen} \textsuperscript{(20)} zu gehen. Nach dem
Essen gingen wir in das \VUgras{Atelier} des Künstlers, wo wir ziemlich lange warten
mussten. Um uns die Zeit zu vertreiben, sahen wir uns die \VUgras{Bildnisse}
\textsuperscript{(22)} an der Wand an.

%%K t14-07-2 p
Als wir an der Reihe waren, ging es schnell, und sobald mein Freund vor der
\VUgras{Kamera} \textsuperscript{(21)} fertig posiert hatte, gingen wir wieder
hinunter.

%%K t14-08-1 p
8. --- Im ersten Stock sahen wir durch die offene Tür des \VUgras{Friseurs}
\textsuperscript{(23)}, wie sein Gehilfe einem Kunden die Haare schnitt.

%%K t14-08-2 p
Wieder auf der Straße, kamen wir am \VUgras{Tabakladen} \textsuperscript{(24)}
vorbei. Johann amüsierte sich so über die rote \VUgras{Laterne}
\textsuperscript{(25)} und die \VUgras{große Pfeife}, die als \VUgras{Schild}
\textsuperscript{(26)} dient, dass er in zwei alte Herren hineinlief, die sich bei
einer \VUgras{Prise Schnupftabak} \textsuperscript{(27)} unterhielten.

%%K t14-tit-4 sub
\VUpk{10.2pt}{40}{Die Konditorei.}
\VUsaut{3.0mm}

%%K t14-09-1 p
9. --- Die \VUgras{Geldscheine} und \VUgras{Münzen} aus allen Ländern im
\VUgras{Schaufenster} des \VUgras{Geldwechslers} \textsuperscript{(29)} hielten uns
einen Augenblick fest, aber weil es Zeit zum Kaffee war, gingen wir ohne weiteren
Aufenthalt in die \VUgras{Konditorei} \textsuperscript{(31)}.

%%K t14-10-1 p
10. --- Beim Anblick all der \VUgras{guten Sachen} im Laden --- \VUgras{Torten},
\VUgras{Lebkuchen}, \VUgras{Gebäck} und \VUgras{Süßigkeiten} jeder Art --- konnten
wir uns kaum entscheiden. Am Ende nahm Johann die \VUgras{Sahnetörtchen}, und ich
nahm nur ein kleines \VUgras{Kirschtörtchen}, weil ich von Süßem
\VUgras{Zahnschmerzen bekomme} und lieber nicht zu oft zum \VUgras{Zahnarzt}
\textsuperscript{(30)} gehe.

%%K t14-tit-5 sub
\VUpk{10.2pt}{40}{Maschinenschreiben.}
\VUsaut{3.0mm}

%%K t14-11-1 p
11. --- Um meinem Freund eine \VUgras{Schreibmaschine} zu zeigen, von der er gehört,
die er aber nie gesehen hatte, gingen wir in das \VUgras{Büro} \textsuperscript{(32)}
meines \VUgras{Vetters} \textsuperscript{(34)}. Wir fanden ihn dabei, seiner
\VUgras{Sekretärin} \textsuperscript{(35)}, die \VUgras{Stenografin} ist, Briefe zu
diktieren. Kaum war ein Brief fertig, tippte ihn eine \VUgras{Maschinenschreiberin}
\textsuperscript{(33)} auf ihrer Maschine ab. Wir wollten die Arbeit meines
Verwandten nicht stören und blieben nur ein paar Minuten.

%%K t14-12-1 p
12. --- Unter dem Büro meines Vetters wohnt ein großer
\VUgras{Uhrmacher und Juwelier} \textsuperscript{(37)}, der auch Silberschmied ist.
Sein prächtiger Laden hat viele \VUgras{Uhren}, \VUgras{Kaminuhren},
\VUgras{Taschenuhren}, \VUgras{Ketten} und kostbare \VUgras{Juwelen} wie
\VUgras{Perlenketten}, goldene \VUgras{Armbänder}, \VUgras{Ohrringe} und
\VUgras{Ringe} mit \VUgras{Edelsteinen} und \VUgras{Diamanten}. Eines der
Schaufenster ist ganz dem \VUgras{Silbergerät} gewidmet und enthält eine große
Sammlung silbernen \VUgras{Bestecks}.

%%K t14-13-1 p
13. --- Als wir um die Ecke der Straße bogen, fiel mir auf, dass das
\VUgras{Namensschild} \textsuperscript{(36)} neben der \VUgras{Hausnummer}
gewechselt worden war. Ich wollte es gerade laut sagen, als wir den fröhlichen Klang
einer Militärkapelle hörten. Wir liefen dem Regiment entgegen, ohne daran zu denken,
dass es an den Läden des \VUgras{Hutmachers} \textsuperscript{(39)} und des
\VUgras{Handschuhmachers} \textsuperscript{(40)} vorbeikommen würde und dass wir es
ebenso gut hätten vorbeimarschieren sehen können, wenn wir vor dem Schaufenster des
\VUgras{Optikers} \textsuperscript{(38)} geblieben wären, voller \VUgras{Zwicker},
\VUgras{Brillen} und \VUgras{Operngläser}.

%%K t14-tit-6 sub
\VUpk{10.2pt}{40}{Der Vorbeimarsch.}
\VUsaut{3.0mm}

%%K t14-14-1 p
14. --- An der Spitze des \VUgras{Regiments} \textsuperscript{(41)} marschierte der
\VUgras{Tambourmajor} \textsuperscript{(42)}, hinter ihm die \VUgras{Trommler}
\textsuperscript{(43)}, die \VUgras{Hornisten} \textsuperscript{(44)} und die ganze
\VUgras{Kapelle}. Der \VUgras{General} \textsuperscript{(45)}, der mit seinem
Adjutanten auf dem Platz war, hatte neben dem \VUgras{Wasserstrahl}
\textsuperscript{(49)} des \VUgras{Springbrunnens} \textsuperscript{(48)} angehalten,
um die \VUgras{Parade} anzusehen.

%%K t14-14-2 p
\VUgras{Die Stabsoffiziere} \textsuperscript{(46)}, der \VUgras{Oberst} und der
\VUgras{Oberstleutnant} grüßten ihn mit dem Degen. \VUgras{Die Majore} standen an
der Spitze ihrer \VUgras{Bataillone} \textsuperscript{(47)}, und jeder
\VUgras{Hauptmann} führte seine \VUgras{Kompanie}, während die \VUgras{Leutnants},
die \VUgras{Fähnriche} und die \VUgras{Unteroffiziere} ihre gewohnten Plätze neben
ihren Leuten einnahmen.

%%K t14-15-1 p
15. --- Viele Vorübergehende hatten sich an der \VUgras{Kreuzung}
\textsuperscript{(50)} versammelt; die Ladenbesitzer --- der \VUgras{Schirmhändler}
\textsuperscript{(51)}, der \VUgras{Färber} \textsuperscript{(53)}, der
\VUgras{Büchsenmacher} \textsuperscript{(54)} --- kamen heraus, um die
\VUgras{Soldaten} zu sehen. Alle Fahrzeuge --- \VUgras{Droschken}
\textsuperscript{(52)}, \VUgras{Privatwagen} \textsuperscript{(57)} und auch der
\VUgras{Sprengwagen} \textsuperscript{(56)} --- fuhren an den Bürgersteig heran.

%%K t14-tit-7 sub
\VUpk{10.2pt}{40}{Bilder von der Straße.}
\VUsaut{3.0mm}

%%K t14-15-2 p
Ein \VUgras{Handelsreisender} \textsuperscript{(55)} mit seinen
\VUgras{Musterkoffern} in der Hand überquerte rasch die Straße, und die Leute auf
dem \VUgras{Oberdeck} \textsuperscript{(59)} des \VUgras{Omnibusses}
\textsuperscript{(58)} drehten sich um, um das Schauspiel zu genießen. Ein armer
\VUgras{Straßensänger} \textsuperscript{(60)} mit seiner \VUgras{Drehorgel}
\textsuperscript{(61)} musste sein \VUgras{Lied} abbrechen, weil der Lärm der Trommeln
seine Stimme übertönte.

%%K t14-16-1 p
16. --- Als das Regiment vor das \VUgras{Tuchgeschäft} \textsuperscript{(64)} kam,
ließen die \VUgras{Näherinnen} \textsuperscript{(63)} und die \VUgras{Schneider}
\textsuperscript{(62)} ihre \VUgras{Nähmaschinen} und ihre Arbeit stehen und sahen
zum Fenster hinaus. \VUgras{Der Ladenbesitzer} \textsuperscript{(65)}, der eben neue
\VUgras{Anzüge} \textsuperscript{(66)} in seine \VUgras{Auslage} gehängt hatte, musste
die Ballen \VUgras{Tuch} \textsuperscript{(67)} und \VUgras{Leinen}, die auf dem
Bürgersteig neben dem \VUgras{Gullyloch} \textsuperscript{(68)} lagen, wegräumen, aus
Furcht, die Menge könnte sie umwerfen; selbst ein alter \VUgras{Hausierer}
\textsuperscript{(69)}, mit dem eine Hausmeisterin gerade um ein paar Strümpfe
feilschte, musste in einen Hausflur treten, um sein Geschäft zu Ende zu bringen.

%%K t14-16-2 p
Wir begleiteten das Regiment bis zur Kaserne \VUgras{und} kamen sehr zufrieden mit
unserem Tag zum Abendessen nach Hause.

%%K t14-tit-8 sub
\VUpk{10.2pt}{40}{Allerlei Gewerbe und Berufe.}
\VUsaut{3.0mm}

%%K t14-17-1 p
17. --- Am nächsten Tag bot mein Vater an, uns die Druckerei der Lokalzeitung zu
zeigen. Wir sagten begeistert zu.

%%K t14-17-2 p
Der \VUgras{Drucker} ist zugleich \VUgras{Buchhändler} \textsuperscript{(70)},
\VUgras{Verleger}, \VUgras{Papierhändler} und \VUgras{Buchbinder}. Er hat in den
ersten Stock die \VUgras{Redaktion} \textsuperscript{(71)} der Zeitung gelegt und auch
den \VUgras{Stichraum}, wo die \VUgras{Lithografen} \textsuperscript{(72)} und die
\VUgras{Kupferstecher} \textsuperscript{(73)} arbeiten, und schließlich den Setzsaal,
wo die \VUgras{Setzer} \textsuperscript{(74)} sind. Der eigentliche
\VUgras{Maschinensaal} \textsuperscript{(75)}, die \VUgras{Druckpressen} und die
verschiedenen Maschinen sind im Erdgeschoss.

%%K t14-tit-9 sub
\VUpk{10.2pt}{40}{Johann fragt sehr viel.}
\VUsaut{3.0mm}

%%K t14-18-1 p
18. --- Als wir aus der Druckerei kamen, blieb Johann ganz verwundert vor einem
kleinen Glaskasten auf einer Säule gegenüber dem \VUgras{Kurzwarenladen}
\textsuperscript{(77)} stehen und fragte, wozu er da sei. Mein Vater erklärte ihm,
das sei ein \VUgras{Feuermelder} \textsuperscript{(76)}. Überhaupt war meinem Freund
in der Stadt alles ein Wunder, vom schlichten \VUgras{Trinkbrunnen}
\textsuperscript{(78)} und dem leichten \VUgras{Lieferdreirad}
\textsuperscript{(80)}, das ein \VUgras{Laufbursche} \textsuperscript{(79)} in Livree
fuhr, bis zum \VUgras{Postwagen} \textsuperscript{(81)}.

%%K t14-19-1 p
19. --- Während mein Vater uns einen Augenblick verließ, um mit dem
\VUgras{Steuereinnehmer} \textsuperscript{(83)} zu sprechen, der stehen geblieben war,
um mit einem \VUgras{Rechtsanwalt} \textsuperscript{(82)} und einem \VUgras{Notar}
\textsuperscript{(84)} zu reden, vertrieben wir uns die Zeit damit, dem Verkehr auf
der Straße zu folgen. Die verschiedensten Leute zogen an uns vorbei: ein
\VUgras{Konditorlehrling} \textsuperscript{(85)}, der eine prächtige \VUgras{Pastete}
in einem Korb trug, kam hinter einem \VUgras{Kleiderhändler}
\textsuperscript{(86)}; ein \VUgras{Laternenanzünder} \textsuperscript{(87)}
unterhielt sich mit einem \VUgras{Gasinstallateur} \textsuperscript{(88)}; eine
\VUgras{Nonne} \textsuperscript{(89)}, die ein kleines Waisenmädchen an der Hand
hielt, ging in ihr Kloster zurück.

%%K t14-tit-10 sub
\VUpk{10.2pt}{40}{Auf dem Heimweg.}
\VUsaut{3.0mm}

%%K t14-20-1 p
20. --- Gerade als mein Vater wieder zu uns stieß, brach Johann in Lachen aus über
die rußigen Gesichter und die seltsame Aufmachung zweier \VUgras{Schornsteinfeger}
\textsuperscript{(91)}, die schwarz vom Ruß waren.

%%K t14-21-1 p
21. --- Ich wollte meiner Mutter bei der \VUgras{Blumenhändlerin}
\textsuperscript{(92)} eine Pflanze kaufen, aber mein Vater wies mich darauf hin, dass
sie sehr schwer zu tragen wäre und dass \VUgras{Apfelsinen} besser seien. Wir gingen
zur \VUgras{Standbesitzerin} \textsuperscript{(94)}, und als die \VUgras{Erzieherin}
\textsuperscript{(93)}, die für ihren Schützling welche aussuchte, fertig eingekauft
hatte, wurden wir bedient.

%%K t14-22-1 p
22. --- Auf dem Heimweg trafen wir mehrere Bekannte: eine alte Freundin meiner
Familie, die sich auf den Arm ihrer \VUgras{Gesellschafterin} \textsuperscript{(95)}
stützte, den \VUgras{Ingenieur} \textsuperscript{(96)} für Straßen- und Brückenbau und
eine meiner Kusinen mit ihrer \VUgras{Kinderfrau} \textsuperscript{(98)}.

%%K t14-22-2 p
Dann kamen wir an der \VUgras{Putzmacherin} \textsuperscript{(97)} meiner Mutter
vorbei und an zwei \VUgras{Mönchen} \textsuperscript{(99)}, die fremd in der Stadt
waren und meinen Vater nach dem Weg zum Bahnhof fragten.
\VUsaut{6.0mm}
\VUfilet{22mm}
